% Section 3.10, from lectures/L03/README.md: the objectives, the questions to test yourself, and
% the next lecture.

\needspace{10\baselineskip}
\section{Review}
\label{c3:sec:review}

\subsection*{What you should be able to do now}
After this chapter, you should be able to:
\begin{itemize}
  \item Say why sequential logic exists, explain the functional difference between a D latch and a
    D flip-flop, and why synchronous designs are built almost entirely from the flip-flop.
  \item Read and write the synchronous process template, and recognize a register as several
    flip-flops sharing a clock, reset and enable.
  \item Say when a signal assignment takes effect: \code{<=} schedules rather than applies, so two
    assignments in one clocked process build a chain, and their order makes no difference.
  \item Build an edge detector from a flip-flop and a gate, and say exactly what signal it produces
    and for how long.
\end{itemize}

\needspace{6\baselineskip}
\subsection*{Questions to test yourself}
You should be able to explain:
\begin{enumerate}
  \item What is the functional difference between a D latch and a D flip-flop, and under what
    condition does each change its output?
  \item Why do synchronous circuits update their state on a single clock edge instead of
    continuously reacting to their inputs?
  \item Given a signal and a flip-flop holding its value from one cycle earlier, how do you combine
    them into a pulse that is high for exactly one clock cycle when the signal rises?
  \item Why must \code{clock} be the only signal in a flip-flop process's sensitivity list, unless
    the design has an asynchronous reset?
  \item A clocked process contains \code{b <= a;} and then \code{c <= b;}. After one clock edge, what
    does \code{c} hold, and why is it not \code{a}? What changes if you swap the two lines, and why?
\end{enumerate}

\needspace{6\baselineskip}
\subsection*{Looking ahead}
Chapter~\ref{c4:ch} is about metastability and synchronization:
\begin{itemize}
  \item Why the raw button feeding this chapter's edge detector can drive a flip-flop into an
    undefined, unstable state, and why that is a different kind of problem from anything here.
  \item The double-flop synchronizer, and what it does and does not do about a bouncing button.
  \item The same LED toggle, rebuilt so it is actually safe on real hardware.
\end{itemize}
